\documentclass[oneside]{article}

\usepackage{bibentry}
\nobibliography*

\usepackage{geometry}

\usepackage[colorlinks,urlcolor=blue,linkcolor=blue,citecolor=blue]{hyperref}

\begin{filecontents}{\jobname.bib}
@techreport{stratification-undecidable,
    title       = {Undecidability of Stratification of Commutativity Relations},
    author      = {Farzan, Azadeh},
    institution = {University of Toronto},
    year        = {2026},
    note        = {Available at \url{https://www.cs.toronto.edu/~azadeh/resources/papers/popl23-errata.pdf}}
}
\end{filecontents}

%opening
\title{Errata for the POPL'23 paper\\ \em``Stratified Commutativity in Verification Algorithms for Concurrent Programs''}
\author{Dominik Klumpp}

\begin{document}
\maketitle

\section*{Publication}
\bibentry{popl23:stratifiedCommutativity}

\section*{Errata}

\begin{itemize}
\item Section~6 of the paper presents an algorithm that checks whether a given proof covers every interleaving of a concurrent program, modulo \emph{stratified commutativity}.
  It is claimed (Theorem~6.12) that the algorithm precisely decides the premise of the so-called \emph{decidable $n$-stratified proof rule} -- inclusion~(4) in the paper.

  In fact, only one implication direction holds (\emph{soundness}):
  if the algorithm claims the inclusion holds, then it does indeed hold.
  The algorithm is however not complete.

  Section~4.4 of my thesis~\cite{diss:commutativityToProofsAndBack} presents a corrected picture,
  and also clarifies that the algorithm nevertheless generalizes the proof rules based on a single commutativity relation (Theorem~4.36).
%
  The question of whether the premise of the so-called \emph{decidable $n$-stratified proof rule} is indeed (precisely) decidable has since been answered in the negative~\cite{stratification-undecidable}.
\end{itemize}


\bibliographystyle{acm}
\bibliography{\jobname,../../_site/resources/bibliography}

\end{document}
